% GENERATED FILE. Edit canonical lessons, then run npm run generate:books.
% canonical-source-hash: fnv1a64:0bd10e2b1d429a67
% canonical-lessons: TE-C216-snanapu-gadi, TE-C216-punadi, TE-C216-getu, TE-C216-bekari, TE-C216-selunu, TE-R216-first-pass-words-from-chapters-208-211, TE-R216-second-pass-words-from-chapters-212-216

\chapter{Bathroom, Foundation, Gate, Bakery, Hair salon}
\label{ch:te-bathroom-foundation-gate-bakery-hair-salon}
\hlchaptermodality{216}

\begin{chapteropening}
\textbf{By the end of this chapter:} \emph{I can say a bathroom, a foundation, a gate, a bakery and a hair salon.}

Complete the last lesson of chapter 216: I can say a bathroom, a foundation, a gate, a bakery and a hair salon.
\end{chapteropening}

\section[\texorpdfstring{\te{స్నానపు} \te{గది} (snānapu gadi)}{స్నానపు గది (snānapu gadi)}]{\te{స్నానపు} \te{గది} (snānapu gadi) — a bathroom}

\label{lesson:TE-C216-snanapu-gadi}



\begin{quote}
Before the new one: say the Telugu for a theatre, then the Telugu for a stadium.
\end{quote}

\subsection*{You'll want to know: \te{స్నానపు} \te{గది}}
\textbf{\te{స్నానపు} \te{గది}} — \emph{snānapu gadi} — \textquotedblleft{}a bathroom\textquotedblright{}.

\begin{grammarlens}[title={What you've built}]
One more word for this chapter.
\end{grammarlens}

\subsection*{Your turn}
\begin{itemize}
\raggedright
  \item \emph{Say it:} \emph{snānapu gadi}
  \item \emph{Say it:} \emph{snānapu gadi}, once more
  \item \emph{Say it:} say \emph{snānapu gadi}
  \item \emph{Recall:} say the Telugu for a theatre, then the Telugu for a stadium, then say \emph{snānapu gadi} again
\end{itemize}

\begin{tcolorbox}[breakable,colback=teal!4,colframe=teal!35!black,title={Before you move on}]
What does \te{స్నానపు} \te{గది} mean? (\textquotedblleft{}A bathroom\textquotedblright{}.) Say it once more.
\end{tcolorbox}
\section[\texorpdfstring{\te{పునాది} (punādi)}{పునాది (punādi)}]{\te{పునాది} (punādi) — a foundation}

\label{lesson:TE-C216-punadi}



\begin{quote}
Before the new one: say the Telugu for a stadium, then the Telugu for a bathroom.
\end{quote}

\subsection*{You'll want to know: \te{పునాది}}
\textbf{\te{పునాది}} — \emph{punādi} — \textquotedblleft{}a foundation\textquotedblright{}.

\begin{grammarlens}[title={What you've built}]
One more word for this chapter.
\end{grammarlens}

\subsection*{Your turn}
\begin{itemize}
\raggedright
  \item \emph{Say it:} \emph{punādi}
  \item \emph{Say it:} \emph{punādi}, once more
  \item \emph{Say it:} say \emph{punādi}
  \item \emph{Recall:} say the Telugu for a stadium, then the Telugu for a bathroom, then say \emph{punādi} again
\end{itemize}

\begin{tcolorbox}[breakable,colback=teal!4,colframe=teal!35!black,title={Before you move on}]
What does \te{పునాది} mean? (\textquotedblleft{}A foundation\textquotedblright{}.) Say it once more.
\end{tcolorbox}
\section[\texorpdfstring{\te{గేటు} (gēṭu)}{గేటు (gēṭu)}]{\te{గేటు} (gēṭu) — a gate}

\label{lesson:TE-C216-getu}



\begin{quote}
Before the new one: say the Telugu for a bathroom, then the Telugu for a foundation.
\end{quote}

\subsection*{You'll want to know: \te{గేటు}}
\textbf{\te{గేటు}} — \emph{gēṭu} — \textquotedblleft{}a gate\textquotedblright{}.

\begin{grammarlens}[title={What you've built}]
One more word for this chapter.
\end{grammarlens}

\subsection*{Your turn}
\begin{itemize}
\raggedright
  \item \emph{Say it:} \emph{gēṭu}
  \item \emph{Say it:} \emph{gēṭu}, once more
  \item \emph{Say it:} say \emph{gēṭu}
  \item \emph{Recall:} say the Telugu for a bathroom, then the Telugu for a foundation, then say \emph{gēṭu} again
\end{itemize}

\begin{tcolorbox}[breakable,colback=teal!4,colframe=teal!35!black,title={Before you move on}]
What does \te{గేటు} mean? (\textquotedblleft{}A gate\textquotedblright{}.) Say it once more.
\end{tcolorbox}
\section[\texorpdfstring{\te{బేకరీ} (bēkarī)}{బేకరీ (bēkarī)}]{\te{బేకరీ} (bēkarī) — a bakery}

\label{lesson:TE-C216-bekari}



\begin{quote}
Before the new one: say the Telugu for a foundation, then the Telugu for a gate.
\end{quote}

\subsection*{You'll want to know: \te{బేకరీ}}
\textbf{\te{బేకరీ}} — \emph{bēkarī} — \textquotedblleft{}a bakery\textquotedblright{}.

\begin{grammarlens}[title={What you've built}]
One more word for this chapter.
\end{grammarlens}

\subsection*{Your turn}
\begin{itemize}
\raggedright
  \item \emph{Say it:} \emph{bēkarī}
  \item \emph{Say it:} \emph{bēkarī}, once more
  \item \emph{Say it:} say \emph{bēkarī}
  \item \emph{Recall:} say the Telugu for a foundation, then the Telugu for a gate, then say \emph{bēkarī} again
\end{itemize}

\begin{tcolorbox}[breakable,colback=teal!4,colframe=teal!35!black,title={Before you move on}]
What does \te{బేకరీ} mean? (\textquotedblleft{}A bakery\textquotedblright{}.) Say it once more.
\end{tcolorbox}
\section[\texorpdfstring{\te{సెలూను} (selūnu)}{సెలూను (selūnu)}]{\te{సెలూను} (selūnu) — a hair salon}

\label{lesson:TE-C216-selunu}



\begin{quote}
Before the new one: say the Telugu for a gate, then the Telugu for a bakery.
\end{quote}

\subsection*{You'll want to know: \te{సెలూను}}
\textbf{\te{సెలూను}} — \emph{selūnu} — \textquotedblleft{}a hair salon\textquotedblright{}.

\begin{grammarlens}[title={What you've built}]
One more word for this chapter.
\end{grammarlens}

\subsection*{Your turn}
\begin{itemize}
\raggedright
  \item \emph{Say it:} \emph{selūnu}
  \item \emph{Say it:} \emph{selūnu}, once more
  \item \emph{Say it:} say \emph{selūnu}
  \item \emph{Recall:} say the Telugu for a gate, then the Telugu for a bakery, then say \emph{selūnu} again
\end{itemize}

\begin{tcolorbox}[breakable,colback=teal!4,colframe=teal!35!black,title={Before you move on}]
What does \te{సెలూను} mean? (\textquotedblleft{}A hair salon\textquotedblright{}.) Say it once more.
\end{tcolorbox}
\section[(dialogue)]{Words from chapters 208-211 — a first pass}

\label{lesson:TE-R216-first-pass-words-from-chapters-208-211}



\begin{quote}
Say the words for a bakery and a hair salon.
\end{quote}

\subsection*{Your turn}
Say these aloud:

\begin{itemize}
\raggedright
  \item to stick, to gargle, to die, to be born, to confront
  \item an itch, a nurse, a test, medical treatment, the forehead
  \item doubt, contentment, disappointment, boredom, pining
  \item numbness, hiccups, snoring, sunstroke, steam
\end{itemize}

\begin{tcolorbox}[breakable,colback=teal!4,colframe=teal!35!black,title={Before you move on}]
To stick? (\textbf{\te{అతుకు}}, \emph{atuku}.) To gargle? (\textbf{\te{పుక్కిలించు}}, \emph{pukkiliñcu}.) To die? (\textbf{\te{చనిపోవు}}, \emph{canipōvu}.) To be born? (\textbf{\te{పుట్టు}}, \emph{puṭṭu}.) To confront? (\textbf{\te{ఎదుర్కొను}}, \emph{edurkonu}.) An itch? (\textbf{\te{దురద}}, \emph{durada}.) A nurse? (\textbf{\te{నర్సు}}, \emph{narsu}.) A test? (\textbf{\te{పరీక్ష}}, \emph{parīkṣa}.) Medical treatment? (\textbf{\te{వైద్యం}}, \emph{vaidyaṁ}.) The forehead? (\textbf{\te{నుదురు}}, \emph{nuduru}.) Doubt? (\textbf{\te{అనుమానం}}, \emph{anumānaṁ}.) Contentment? (\textbf{\te{తృప్తి}}, \emph{tṛpti}.) Disappointment? (\textbf{\te{నిరాశ}}, \emph{nirāśa}.) Boredom? (\textbf{\te{విసుగు}}, \emph{visugu}.) Pining? (\textbf{\te{బెంగ}}, \emph{beṅga}.) Numbness? (\textbf{\te{తిమ్మిరి}}, \emph{timmiri}.) Hiccups? (\textbf{\te{ఎక్కిళ్ళు}}, \emph{ekkiḷḷu}.) Snoring? (\textbf{\te{గురక}}, \emph{guraka}.) Sunstroke? (\textbf{\te{వడదెబ్బ}}, \emph{vaḍadebba}.) Steam? (\textbf{\te{ఆవిరి}}, \emph{āviri}.)
\end{tcolorbox}
\section[(dialogue)]{Words from chapters 212-216 — a second pass}

\label{lesson:TE-R216-second-pass-words-from-chapters-212-216}



\begin{quote}
Say the words for free time and immediately.
\end{quote}

\subsection*{Your turn}
Say these aloud:

\begin{itemize}
\raggedright
  \item free time, immediately, finally, until, Sankranti
  \item suddenly, again and again, a little while, a second, the day before yesterday
  \item a bus stand, an airport, a junction, a capital city, a district
  \item a desert, a rock, an open ground, a theatre, a stadium
  \item a bathroom, a foundation, a gate, a bakery, a hair salon
\end{itemize}

\begin{tcolorbox}[breakable,colback=teal!4,colframe=teal!35!black,title={Before you move on}]
Free time? (\textbf{\te{తీరిక}}, \emph{tīrika}.) Immediately? (\textbf{\te{వెంటనే}}, \emph{veṇṭanē}.) Finally? (\textbf{\te{చివరకు}}, \emph{civaraku}.) Until? (\textbf{\te{వరకు}}, \emph{varaku}.) Sankranti? (\textbf{\te{సంక్రాంతి}}, \emph{saṅkrānti}.) Suddenly? (\textbf{\te{అకస్మాత్తుగా}}, \emph{akasmāttugā}.) Again and again? (\textbf{\te{పదే} \te{పదే}}, \emph{padē padē}.) A little while? (\textbf{\te{కాసేపు}}, \emph{kāsēpu}.) A second? (\textbf{\te{సెకను}}, \emph{sekanu}.) The day before yesterday? (\textbf{\te{మొన్న}}, \emph{monna}.) A bus stand? (\textbf{\te{బస్టాండు}}, \emph{basṭāṇḍu}.) An airport? (\textbf{\te{విమానాశ్రయం}}, \emph{vimānāśrayaṁ}.) A junction? (\textbf{\te{కూడలి}}, \emph{kūḍali}.) A capital city? (\textbf{\te{రాజధాని}}, \emph{rājadhāni}.) A district? (\textbf{\te{జిల్లా}}, \emph{jillā}.) A desert? (\textbf{\te{ఎడారి}}, \emph{eḍāri}.) A rock? (\textbf{\te{బండ}}, \emph{baṇḍa}.) An open ground? (\textbf{\te{మైదానం}}, \emph{maidānaṁ}.) A theatre? (\textbf{\te{థియేటరు}}, \emph{thiyēṭaru}.) A stadium? (\textbf{\te{స్టేడియం}}, \emph{sṭēḍiyaṁ}.) A bathroom? (\textbf{\te{స్నానపు} \te{గది}}, \emph{snānapu gadi}.) A foundation? (\textbf{\te{పునాది}}, \emph{punādi}.) A gate? (\textbf{\te{గేటు}}, \emph{gēṭu}.) A bakery? (\textbf{\te{బేకరీ}}, \emph{bēkarī}.) A hair salon? (\textbf{\te{సెలూను}}, \emph{selūnu}.)
\end{tcolorbox}
